\documentclass[a4paper,10pt]{article}

% Настройки русского языка и кодировок
\usepackage[T2A]{fontenc}
\usepackage[utf8]{inputenc}
\usepackage[russian]{babel}

% Математические пакеты
\usepackage{amsmath, amsfonts, amssymb, amsthm}
\usepackage{geometry}
\usepackage{titlesec}

% Настройка полей (компактные поля для конспекта)
\geometry{left=1cm, right=1cm, top=1.5cm, bottom=1.5cm}

% Настройка заголовков
\titleformat{\section}{\large\bfseries}{}{0em}{}
\titleformat{\subsection}{\normalsize\bfseries}{}{0em}{}

\begin{document}

\begin{center}
    \LARGE \textbf{Методы решения многомерного волнового уравнения}
\end{center}

\section{1. Постановка задачи Коши для волнового уравнения}
Рассматривается задача Коши для функции $u(x, t) \in C^2(\mathbb{R}^n \times \{t > 0\}) \cap C^1(\mathbb{R}^n \times \{t \ge 0\})$:
\[
\begin{cases}
u_{tt} = a^2 \Delta u + f(x,t) & (1) \\
u|_{t=0} = u_0(x) & (2) \\
u_t|_{t=0} = u_1(x) & (3)
\end{cases}
\]

\section{2. Общие аналитические формулы}
\textbf{Теорема:} Если $u_0(x) \in C^3(\mathbb{R}^n)$, $u_1(x) \in C^2(\mathbb{R}^n)$, $f(x,t) \in C^1(\mathbb{R}^n \times[0, +\infty))$, то классическое решение существует и единственно. \\

\noindent\textbf{Для $n=2$, $x \in \mathbb{R}^2, t>0$ (Формула Пуассона):}
\[ u = \frac{1}{2\pi a} \iint\limits_{|\xi-x|<at} \frac{u_1(\xi)d\xi}{\sqrt{a^2t^2-|\xi-x|^2}} + \frac{1}{2\pi a}\frac{\partial}{\partial t}\left( \iint\limits_{|\xi-x|<at} \frac{u_0(\xi)d\xi}{\sqrt{a^2t^2-|\xi-x|^2}} \right) + \frac{1}{2\pi a}\int\limits_0^t d\tau \iint\limits_{|\xi-x|<a(t-\tau)} \frac{f(\xi, \tau)d\xi}{\sqrt{a^2(t-\tau)^2 - |\xi-x|^2}} \]

\noindent\textbf{Для $n=3$, $x \in \mathbb{R}^3, t>0$ (Формула Кирхгофа):}
\[ u = \frac{1}{4\pi a^2 t} \iint\limits_{|\xi-x|=at} u_1(\xi)dS_\xi + \frac{\partial}{\partial t}\left( \frac{1}{4\pi a^2 t} \iint\limits_{|\xi-x|=at} u_0(\xi)dS_\xi \right) + \frac{1}{4\pi a^2}\iiint\limits_{|\xi-x|<at} \frac{1}{|\xi-x|} f\left(\xi, t - \frac{|\xi-x|}{a}\right) d\xi \]

\textit{Примечание:} На практике прямого вычисления данных интегралов стараются избегать, применяя специальные методы понижения размерности или упрощения задачи (см. ниже).

\section{3. ПРИЁМ 1: Разбивка на сумму задач}
\textbf{Суть:} В силу линейности дифференциального оператора решение исходной задачи ищется в виде суммы $u = v + w + \tilde{u}$, где каждая функция отвечает только за одно ненулевое начальное условие:

\vspace{0.2cm}
\noindent
\begin{minipage}{0.32\textwidth}
\textbf{A)} $v$ (отвечает за $u_0$): 
\[ \begin{cases} v_{tt} = a^2\Delta v \\ v|_{t=0} = u_0(x) \\ v_t|_{t=0} = 0 \end{cases} \]
\end{minipage}
\hfill
\begin{minipage}{0.32\textwidth}
\textbf{B)} $w$ (отвечает за $u_1$): 
\[ \begin{cases} w_{tt} = a^2\Delta w \\ w|_{t=0} = 0 \\ w_t|_{t=0} = u_1(x) \end{cases} \]
\end{minipage}
\hfill
\begin{minipage}{0.32\textwidth}
\textbf{C)} $\tilde{u}$ (отвечает за $f(x,t)$): 
\[ \begin{cases} \tilde{u}_{tt} = a^2\Delta \tilde{u} + f \\ \tilde{u}|_{t=0} = 0 \\ \tilde{u}_t|_{t=0} = 0 \end{cases} \]
\end{minipage}

\section{4. ПРИЁМ 2: Отщепление гармонического множителя}
\textbf{Суть:} Понижение размерности задачи за счет отделения гармонической части. Применяется, если условия задачи имеют вид: 
\[
\begin{cases} 
u_{tt} = a^2 \Delta u + g(x_1, \dots, x_m, t) \cdot h(x_{m+1}, \dots, x_n) \\ 
u|_{t=0} = u_0(x_1, \dots, x_m) \cdot h(x_{m+1}, \dots, x_n) \\ 
u_t|_{t=0} = u_1(x_1, \dots, x_m) \cdot h(x_{m+1}, \dots, x_n) 
\end{cases}
\]
Причем $\Delta h = 0$ (функция $h$ является гармонической по своим переменным).

\noindent\textbf{Решение:} Вспомогательная функция $v = v(x_1, \dots, x_m, t)$ ищется из задачи меньшей размерности: 
\[
\begin{cases} 
v_{tt} = a^2 \Delta v + g(x_1, \dots, x_m, t) \\ 
v|_{t=0} = u_0(x_1, \dots, x_m) \\ 
v_t|_{t=0} = u_1(x_1, \dots, x_m) 
\end{cases}
\]
Тогда итоговое решение имеет вид: $u(x, t) = v(x_1, \dots, x_m, t) \cdot h(x_{m+1}, \dots, x_n)$.

\section{5. ПРИЁМ 3: Использование собственных функций оператора $\Delta$}
\textbf{Скалярный случай:} Если все заданные функции пропорциональны функции $g(x)$, удовлетворяющей условию $\Delta g(x) = \lambda g(x)$: 
\[
\begin{cases} 
u_{tt} = a^2 \Delta u + f(t)g(x) \\ 
u|_{t=0} = \alpha g(x) \\ 
u_t|_{t=0} = \beta g(x) 
\end{cases} 
\implies \text{Решение ищется в виде } u(x,t) = \varphi(t)g(x).
\]
Функция $\varphi(t)$ находится из решения обыкновенного дифференциального уравнения (ОДУ):
\[
\begin{cases} 
\varphi''(t) - \lambda a^2 \varphi(t) = f(t) \\ 
\varphi(0) = \alpha, \quad \varphi'(0) = \beta 
\end{cases}
\]

\noindent\textbf{Векторный случай (системы):} Если правые части представляют собой линейные комбинации $\{g_1, g_2\}$, где $\begin{pmatrix} \Delta g_1 \\ \Delta g_2 \end{pmatrix} = A \begin{pmatrix} g_1 \\ g_2 \end{pmatrix}$:
\[
\begin{cases}
u_{tt} = a^2 \Delta u + f_1(t)g_1(x) + f_2(t)g_2(x) \\
u|_{t=0} = \alpha_1 g_1(x) + \alpha_2 g_2(x) \\
u_t|_{t=0} = \beta_1 g_1(x) + \beta_2 g_2(x)
\end{cases}
\]
Решение ищется в виде $u(x,t) = \varphi_1(t)g_1(x) + \varphi_2(t)g_2(x)$. Для неизвестных функций возникает система ОДУ: 
\[
\begin{pmatrix} \varphi_1'' \\ \varphi_2'' \end{pmatrix} = 
a^2 A^T \begin{pmatrix} \varphi_1 \\ \varphi_2 \end{pmatrix} + \begin{pmatrix} f_1(t) \\ f_2(t) \end{pmatrix}, \qquad
\begin{pmatrix} \varphi_1(0) \\ \varphi_2(0) \end{pmatrix} = \begin{pmatrix} \alpha_1 \\ \alpha_2 \end{pmatrix}, \quad
\begin{pmatrix} \varphi_1'(0) \\ \varphi_2'(0) \end{pmatrix} = \begin{pmatrix} \beta_1 \\ \beta_2 \end{pmatrix}
\]

\section{6. ПРИЁМ 4: Полиномиальные и гармонические данные}
\textbf{Суть:} Если существует такое $N$, что $\Delta^N(\cdot) \equiv 0$, то решение ищется в виде конечного ряда (разложения в ряд Тейлора по времени). \\
\textit{Замечание: При $N=1$ начальные данные гармонические, что является частным случаем приёма 3 при $\lambda = 0$.}
\begin{enumerate}
    \setlength{\itemsep}{0pt}
    \item Для начального отклонения $u_0$: \quad $u = \sum_{k=0}^{N-1} a^{2k} \frac{t^{2k}}{(2k)!} \Delta^k u_0(x)$ 
    \item Для начальной скорости $u_1$: \quad $u = \sum_{k=0}^{N-1} a^{2k} \frac{t^{2k+1}}{(2k+1)!} \Delta^k u_1(x)$ 
    \item Для правой части $f(x)$: \quad $u = \sum_{k=0}^{N-1} a^{2k} \frac{t^{2k+2}}{(2k+2)!} \Delta^k f(x)$
\end{enumerate}
\textit{Примечание:} Приведенная формула для правой части справедлива только в случае стационарной функции $f = f(x)$. Если $f$ явно зависит от времени ($f = f(x,t)$), данный метод не применяется.

\section{7. ПРИЁМ 5: Сведение к одной переменной (Плоские волны)}
\textbf{Суть:} Если все условия задачи зависят от линейной комбинации переменных $\xi = \alpha_1 x_1 + \dots + \alpha_n x_n = \sum \alpha_i x_i$, причем:
\[ f(x,t) = \tilde{f}(\xi, t)\Big|_{\xi = \sum \alpha_i x_i}, \quad u_0(x) = \tilde{u}_0(\xi)\Big|_{\xi = \sum \alpha_i x_i}, \quad u_1(x) = \tilde{u}_1(\xi)\Big|_{\xi = \sum \alpha_i x_i}. \]
Решение ищется в виде $u(x,t) = v(\xi, t)$. Уравнение сводится к одномерному:
\[ 
\begin{cases}
v_{tt}(\xi,t) = \left( a^2 \sum\limits_{i=1}^n \alpha_i^2 \right) v_{\xi\xi}(\xi,t) + \tilde{f}(\xi, t) \\
v|_{t=0} = \tilde{u}_0(\xi) \\
v_t|_{t=0} = \tilde{u}_1(\xi)
\end{cases}
\]

\section{8. ПРИЁМ 6: Сферическая симметрия ($n=3$)}
Для пространства $\mathbb{R}^3$, если условия зависят только от радиуса $r = |x|$, решение ищется в виде сферически симметричной функции $u(x,t) = v(r, t)$. \\
Вводится замена \textbf{$w(r,t) = r \cdot v(r,t)$}. Для функции $w$ получается одномерная задача на полуоси:
\[ 
\begin{cases} 
w_{tt} = a^2 w_{rr}, \quad r>0, \ t>0 \\ 
w|_{t=0} = r u_0(r), \quad w_t|_{t=0} = r u_1(r) \\ 
w|_{r=0} = 0 
\end{cases} 
\]
Задача решается стандартным методом нечетного продолжения начальных условий на область $r < 0$. 

\vspace{0.2cm}
\noindent\textit{Замечание:} Если требуется найти значение функции в начале координат $u(0,t)$, оно вычисляется по правилу Лопиталя:
\[
u(0,t) = \lim_{r \to 0} \frac{w(r,t)}{r} = w_r(0,t)
\]

\newpage
\appendix
\section{Приложение: Доказательства математических приёмов}

\subsection*{Доказательство к Приёму 1 (Разбивка на сумму задач)}
\begin{proof}
Пусть функции $v, w, \tilde{u}$ являются решениями соответствующих задач A, B и C. Требуется доказать, что их сумма $u(x,t) = v(x,t) + w(x,t) + \tilde{u}(x,t)$ является решением исходной задачи (1)-(3).
В силу линейности оператора дифференцирования подставим сумму в исходное уравнение:
\[
u_{tt} = v_{tt} + w_{tt} + \tilde{u}_{tt} = a^2\Delta v + a^2\Delta w + (a^2\Delta \tilde{u} + f(x,t)) = a^2\Delta(v+w+\tilde{u}) + f(x,t) = a^2\Delta u + f(x,t).
\]
Уравнение выполнено. Проверим начальные условия при $t=0$:
\[
u|_{t=0} = v|_{t=0} + w|_{t=0} + \tilde{u}|_{t=0} = u_0(x) + 0 + 0 = u_0(x).
\]
\[
u_t|_{t=0} = v_t|_{t=0} + w_t|_{t=0} + \tilde{u}_t|_{t=0} = 0 + u_1(x) + 0 = u_1(x).
\]
Все условия задачи Коши удовлетворены. В силу теоремы единственности, найденное $u(x,t)$ есть единственное искомое решение.
\end{proof}

\subsection*{Доказательство к Приёму 2 (Отщепление гармонического множителя)}
\begin{proof}
В общем случае оператор Лапласа от произведения двух функций раскрывается по формуле:
\[ \Delta(v \cdot h) = (\Delta v)h + 2(\nabla v, \nabla h) + v(\Delta h) \]
По условию функция $v$ зависит только от переменных $(x_1, \dots, x_m)$, а функция $h$ — только от $(x_{m+1}, \dots, x_n)$. Следовательно, их градиенты ортогональны, и скалярное произведение градиентов равно нулю: $(\nabla v, \nabla h) = \sum \frac{\partial v}{\partial x_i}\frac{\partial h}{\partial x_i} = 0$. 
Кроме того, по условию функция $h$ является гармонической, то есть $\Delta h = 0$. 
Следовательно, формула упрощается до: $\Delta(v \cdot h) = (\Delta v) \cdot h$. \\
Умножим вспомогательное уравнение для $v$ на функцию $h$:
\[
v_{tt} \cdot h = (a^2\Delta v + g) \cdot h \implies (v \cdot h)_{tt} = a^2(\Delta v) \cdot h + g \cdot h \implies u_{tt} = a^2\Delta u + g \cdot h.
\]
Это в точности совпадает с исходным уравнением. Начальные условия удовлетворяются очевидным образом путем домножения условий для $v$ на функцию $h$.
\end{proof}

\subsection*{Доказательство к Приёму 3 и 3* (Собственные функции оператора $\Delta$)}
\begin{proof}
\textbf{Для векторного случая 3* (скалярный доказывается аналогично при $A=\lambda$):} \\
Ищем решение в виде $u(x,t) = \varphi_1(t)g_1(x) + \varphi_2(t)g_2(x)$. Применим к нему оператор Лапласа. По условию задана матрица $A = \begin{pmatrix} a & b \\ c & d \end{pmatrix}$, такая что $\Delta g_1 = a g_1 + b g_2$ и $\Delta g_2 = c g_1 + d g_2$. Тогда:
\[
\begin{aligned}
\Delta u &= \varphi_1\Delta g_1 + \varphi_2\Delta g_2 = \varphi_1(ag_1 + bg_2) + \varphi_2(cg_1 + dg_2) \\
&= (a\varphi_1 + c\varphi_2)g_1 + (b\varphi_1 + d\varphi_2)g_2.
\end{aligned}
\]
Подставим предполагаемое решение в исходное волновое уравнение $u_{tt} = a^2 \Delta u + f_1 g_1 + f_2 g_2$:
\[
\varphi_1'' g_1 + \varphi_2'' g_2 = a^2 \left[ (a\varphi_1 + c\varphi_2)g_1 + (b\varphi_1 + d\varphi_2)g_2 \right] + f_1 g_1 + f_2 g_2.
\]
Перегруппируем члены при базисных функциях $g_1(x)$ и $g_2(x)$ и приравняем их:
\[
\begin{cases}
\varphi_1'' = a^2(a\varphi_1 + c\varphi_2) + f_1 \\
\varphi_2'' = a^2(b\varphi_1 + d\varphi_2) + f_2
\end{cases}
\]
Заметим, что коэффициенты в правой части образуют матрицу, транспонированную по отношению к $A$. В матричной форме это записывается как:
\[
\begin{pmatrix} \varphi_1'' \\ \varphi_2'' \end{pmatrix} = a^2 \begin{pmatrix} a & c \\ b & d \end{pmatrix} \begin{pmatrix} \varphi_1 \\ \varphi_2 \end{pmatrix} + \begin{pmatrix} f_1 \\ f_2 \end{pmatrix} = a^2 A^T \begin{pmatrix} \varphi_1 \\ \varphi_2 \end{pmatrix} + \begin{pmatrix} f_1 \\ f_2 \end{pmatrix}.
\]
\end{proof}

\subsection*{Доказательство к Приёму 4 (Полиномиальные данные)}
\begin{proof}
Рассмотрим случай, когда задано только начальное отклонение $u_0(x)$ (для остальных пунктов доказательство строится полностью аналогично), причем $\Delta^N u_0(x) \equiv 0$. Подставим ряд $u = \sum_{k=0}^{N-1} a^{2k} \frac{t^{2k}}{(2k)!} \Delta^k u_0(x)$ в уравнение $u_{tt} = a^2\Delta u$. \\
Вычислим вторую производную по времени (при $k=0$ производная константы равна 0, поэтому суммирование начинается с $k=1$):
\[
u_{tt} = \sum_{k=1}^{N-1} a^{2k} \frac{2k(2k-1)t^{2k-2}}{(2k)!} \Delta^k u_0(x) = \sum_{k=1}^{N-1} a^{2k} \frac{t^{2k-2}}{(2k-2)!} \Delta^k u_0(x).
\]
Сделаем сдвиг индекса суммирования $l = k - 1$:
\[
u_{tt} = \sum_{l=0}^{N-2} a^{2l+2} \frac{t^{2l}}{(2l)!} \Delta^{l+1} u_0(x).
\]
Теперь вычислим правую часть уравнения $a^2\Delta u$, внеся оператор Лапласа под знак суммы:
\[
a^2\Delta u = a^2 \sum_{k=0}^{N-1} a^{2k} \frac{t^{2k}}{(2k)!} \Delta^{k+1} u_0(x) = \sum_{k=0}^{N-1} a^{2k+2} \frac{t^{2k}}{(2k)!} \Delta^{k+1} u_0(x).
\]
Сравним $u_{tt}$ и $a^2\Delta u$. Они полностью совпадают для всех индексов от $0$ до $N-2$. В правой части остается одно нескомпенсированное слагаемое при $k = N-1$:
\[
a^{2N} \frac{t^{2N-2}}{(2N-2)!} \Delta^N u_0(x).
\]
Однако по условию $\Delta^N u_0(x) \equiv 0$, следовательно, это слагаемое равно нулю. Таким образом, $u_{tt} = a^2\Delta u$ выполняется тождественно. При подстановке $t=0$ все слагаемые ряда, кроме первого ($k=0$), обнуляются, что дает $u|_{t=0} = u_0(x)$.
\end{proof}

\subsection*{Доказательство к Приёму 5 (Сведение к одной переменной, плоские волны)}
\begin{proof}
Пусть $u(x,t) = v(\xi, t)$, где $\xi = \sum_{i=1}^n \alpha_i x_i$. Вычислим производные по правилу дифференцирования сложной функции:
\[
\frac{\partial u}{\partial x_i} = \frac{\partial v}{\partial \xi} \cdot \frac{\partial \xi}{\partial x_i} = v_\xi \cdot \alpha_i.
\]
Вторая производная:
\[
\frac{\partial^2 u}{\partial x_i^2} = \frac{\partial}{\partial x_i} (v_\xi \cdot \alpha_i) = v_{\xi\xi} \cdot \alpha_i^2.
\]
Оператор Лапласа есть сумма вторых производных:
\[
\Delta u = \sum_{i=1}^n u_{x_i x_i} = \sum_{i=1}^n \left( v_{\xi\xi} \cdot \alpha_i^2 \right) = v_{\xi\xi} \sum_{i=1}^n \alpha_i^2.
\]
Подставляя это в исходное волновое уравнение $u_{tt} = a^2\Delta u + f$, получаем одномерное волновое уравнение для $v$:
\[
v_{tt} = a^2 \left(\sum_{i=1}^n \alpha_i^2 \right) v_{\xi\xi} + \tilde{f}(\xi, t).
\]
При подстановке $t=0$ функции $u(x,0) = u_0(x)$ и $u_t(x,0) = u_1(x)$ очевидным образом принимают вид функций от новой переменной $\xi$, то есть $v(\xi, 0) = \tilde{u}_0(\xi)$ и $v_t(\xi, 0) = \tilde{u}_1(\xi)$.
\end{proof}

\subsection*{Доказательство к Приёму 6 (Сферическая симметрия)}
\begin{proof}
В сферических координатах для функции, зависящей только от радиуса $r = |x|$, оператор Лапласа имеет вид:
\[ \Delta v = v_{rr} + \frac{2}{r}v_r. \]
Сделаем подстановку $w(r,t) = r \cdot v(r,t)$, откуда $v = \frac{w}{r}$. Выразим производные:
\[
v_t = \frac{w_t}{r}, \quad v_{tt} = \frac{w_{tt}}{r}.
\]
Вычислим производные по радиусу $r$:
\[
v_r = \frac{w_r \cdot r - w}{r^2} = \frac{w_r}{r} - \frac{w}{r^2}.
\]
\[
v_{rr} = \frac{\partial}{\partial r} \left( \frac{w_r}{r} - \frac{w}{r^2} \right) = \left( \frac{w_{rr}}{r} - \frac{w_r}{r^2} \right) - \left( \frac{w_r}{r^2} - \frac{2w}{r^3} \right) = \frac{w_{rr}}{r} - \frac{2w_r}{r^2} + \frac{2w}{r^3}.
\]
Подставим найденные производные в выражение для оператора Лапласа:
\[
\Delta v = v_{rr} + \frac{2}{r}v_r = \left( \frac{w_{rr}}{r} - \frac{2w_r}{r^2} + \frac{2w}{r^3} \right) + \frac{2}{r} \left( \frac{w_r}{r} - \frac{w}{r^2} \right) = \frac{w_{rr}}{r} - \frac{2w_r}{r^2} + \frac{2w}{r^3} + \frac{2w_r}{r^2} - \frac{2w}{r^3} = \frac{w_{rr}}{r}.
\]
Подставляя в уравнение $v_{tt} = a^2\Delta v$, получаем:
\[
\frac{w_{tt}}{r} = a^2 \frac{w_{rr}}{r} \implies w_{tt} = a^2 w_{rr} \quad (r>0).
\]
Чтобы исходная функция $v = \frac{w}{r}$ была ограничена в центре (при $r \to 0$), необходимо потребовать выполнение краевого условия $w(0, t) = 0$.
\end{proof}

\end{document}